\documentclass[11pt]{article}
\usepackage[a4paper,margin=18mm]{geometry}
\usepackage[no-math]{fontspec}
\setmainfont{Latin Modern Roman}
\usepackage{amsmath,amssymb,amsthm,xfrac,xspace,hyperref,graphicx,comment}
\excludecomment{DefTralics}
\providecommand{\C}{}
\providecommand{\bibliofont}{\small}
\newcommand{\ie}{i.e.,\xspace}
\newcommand{\eg}{e.g.,\xspace}
\newcommand{\etc}{etc.\xspace}
\newcommand{\cf}{cf.\xspace}
\newcommand{\resp}{resp.\xspace}
\newcommand{\smaller}{\small}
\newcommand{\Subsection}{\subsection}
\newcommand{\Subsubsection}{\subsubsection}
\newcommand{\skpt}{\smallskip}
\emergencystretch=3em
\providecommand{\dpl}{\displaystyle}
\input{author-preamble.tex}
\begin{document}
\begin{center}{\Large Stable item 1932}\end{center}
\paragraph{Source and attribution.}
Clara L. Aldana, Jean-Baptiste Caillau, Pedro Freitas, \emph{Maximal determinants of Schrödinger operators on bounded intervals}, Journal de l’École polytechnique — Mathématiques 7 (2020), publisher version, DOI 10.5802/jep.128. \href{https://jep.centre-mersenne.org/articles/10.5802/jep.128/}{Publisher article}. Authors' text: \href{https://creativecommons.org/licenses/by/4.0/}{CC BY 4.0}.
\paragraph{Editorial problem boundary.}
Prove only Theorem A below for nonnegative real \(V\in L^1[0,1]\), mass at most \(A>0\), and the Dirichlet Schrödinger operator. The unique maximiser is the centred pulse with the printed length and determinant. The entire rough-potential determinant foundation, bounded-control existence, switching and singular-arc arguments, finite-cap limit and complete uniqueness proof are retained. Signed potentials and the \(q>1\) results are surrounding context only.
\paragraph{Difficulty (editorial): H2.}
Pontryagin switching-function analysis excludes interior zero and saturated arcs and proves a singular-pulse structure. Stationary optimizers of bounded-control approximations establish existence despite the noncompact L1 budget, and an unbounded-control maximum principle gives global uniqueness.
\paragraph{Editorial transcription note.}
Original author language, mathematics and ordering are preserved. Publisher front matter is replaced by this independent article wrapper. Bibliography is recovered from matching cached publisher HTML, including its TeX formulas. No new proof or mathematical repair is supplied. Surrounding original results are retained for conservative context and are not extra selected corpus claims.
\clearpage
\input{author-body.tex}
\begin{thebibliography}{99}
\bibitem{agrachev-2004a} A. A. Agrachev \& Y. L. Sachkov - Control theory from the geometric viewpoint , Encyclopaedia of Math. Sciences , vol. 87 , Springer-Verlag, Berlin, 2004
\bibitem{aar} P. Albin, C. L. Aldana \& F. Rochon - “Ricci flow and the determinant of the Laplacian on non-compact surfaces” , Comm. Partial Differential Equations 38 (2013) no. 4, p. 711-749
\bibitem{aursal} E. Aurell \& P. Salomonson - “On functional determinants of Laplacians in polygons and simplicial complexes” , Comm. Math. Phys. 165 (1994) no. 2, p. 233-259
\bibitem{bonnard-2003a} B. Bonnard \& M. Chyba - Singular trajectories and their role in control theory , Math. \& Applications , vol. 40 , Springer-Verlag, Berlin, 2003
\bibitem{BFK} D. Burghelea, L. Friedlander \& T. Kappeler - “On the determinant of elliptic boundary value problems on a line segment” , Proc. Amer. Math. Soc. 123 (1995) no. 10, p. 3027-3038
\bibitem{BGRV} G. Buttazzo, A. Gerolin, B. Ruffini \& B. Velichkov - “Optimal potentials for Schrödinger operators” , J. Éc. polytech. Math. 1 (2014), p. 71-100
\bibitem{cesari-1983a} L. Cesari - Optimization—theory and applications. Problems with ordinary differential equations , Applications of Math. , vol. 17 , Springer-Verlag, New York, 1983
\bibitem{CL} E. A. Coddington \& N. Levinson - Theory of ordinary differential equations , McGraw-Hill Book Company, Inc., New York-Toronto-London, 1955
\bibitem{Edwards} H. M. Edwards - Riemann’s zeta function , Dover Publications, Inc., Mineola, NY, 2001
\bibitem{frei} P. Freitas - “The spectral determinant of the isotropic quantum harmonic oscillator in arbitrary dimensions” , Math. Ann. 372 (2018) no. 3-4, p. 1081-1101
\bibitem{fl} P. Freitas \& J. Lipovsky - “The determinant of one-dimensional polyharmonic operators of arbitrary order” , 2020
\bibitem{geya} I. M. Gelfand \& A. M. Jaglom - “Integration in functional spaces and its applications in quantum physics” , J. Math. Phys. 1 (1960), p. 48-69
\bibitem{harr} E. M. Harrell - “Hamiltonian operators with maximal eigenvalues” , J. Math. Phys. 25 (1984) no. 1, p. 48-51, Erratum: Ibid. 27 (1986), no. 1, p. 419
\bibitem{henr} A. Henrot - Extremum problems for eigenvalues of elliptic operators , Frontiers in Math. , Birkhäuser Verlag, Basel, 2006
\bibitem{Kato} T. Kato - Perturbation theory for linear operators , Classics in Math. , Springer-Verlag, Berlin, 1995
\bibitem{Lesch} M. Lesch - “Determinants of regular singular Sturm-Liouville operators” , Math. Nachr. 194 (1998), p. 139-170
\bibitem{Lesch-Tolk} M. Lesch \& J. Tolksdorf - “On the determinant of one-dimensional elliptic boundary value problems” , Comm. Math. Phys. 193 (1998) no. 3, p. 643-660
\bibitem{Levit-Smilansky} S. Levit \& U. Smilansky - “A theorem on infinite products of eigenvalues of Sturm-Liouville type operators” , Proc. Amer. Math. Soc. 65 (1977) no. 2, p. 299-302
\bibitem{mipl} S. Minakshisundaram \& Å. Pleijel - “Some properties of the eigenfunctions of the Laplace-operator on Riemannian manifolds” , Canad. J. Math. 1 (1949), p. 242-256
\bibitem{okik} K. Okikiolu - “Critical metrics for the determinant of the Laplacian in odd dimensions” , Ann. of Math. (2) 153 (2001) no. 2, p. 471-531
\bibitem{ops} B. Osgood, R. Phillips \& P. Sarnak - “Extremals of determinants of Laplacians” , J. Funct. Anal. 80 (1988) no. 1, p. 148-211
\bibitem{rasi} D. B. Ray \& I. M. Singer - “$R$-torsion and the Laplacian on Riemannian manifolds” , Adv. Math. 7 (1971), p. 145-210
\bibitem{riem} B. Riemann - “Ueber die Anzahl der Primzahlen unter einer gegebenen Grösse” , Monatsber. Berlin. Akad. (1859), p. 671-680, English translation in [9]
\bibitem{Savchuk} A. M. Savchuk - “On the eigenvalues and eigenfunctions of the Sturm-Liouville operator with a singular potential” , Mat. Zametki 69 (2001) no. 2, p. 277-285
\bibitem{SavcShka99} A. M. Savchuk \& A. A. Shkalikov - “The trace formula for Sturm-Liouville operators with singular potentials” , Mat. Zametki 69 (2001) no. 3, p. 427-442
\bibitem{titc} E. C. Titchmarsh - The theory of the Riemann zeta-function , The Clarendon Press, Oxford University Press, New York, 1986, Ed. by D. R. Heath-Brown
\end{thebibliography}
\end{document}
